\section{Weighted boxes and path packing}\label{sec:boxes}

We now count endpoints separated by a fixed positive distance that can
be reached from one starting point by paths of length $O(r)$ with a
prescribed height range. Fix a nonempty compact convex set $E\subset W$
and a scale $r\ge1$.
At heights in $rE$, a generator of weight $\alpha_i$ can
be used with amplitude of order $\exp(r\max_E\alpha_i)$. The boxes below
record the volume generated by these displacements. They will give
matching exponential bounds: many separated endpoints can be reached
while staying within sublinear distance of $rE$, and no such path family
confined to those heights can have more separated endpoints at
the exponential scale. Comparing these bounds for a quasi-isometry
will constrain how it changes height ranges.

The nilpotent group $N$ need not be abelian, and a bracket of nonzero
weight vectors may have weight zero. We therefore retain every bracket
direction. The estimates also retain the kernel of
$q:N\rtimes D\to G$.

Throughout this section $N\ne\{1\}$. Use the homogeneous generators
$X_1,\ldots,X_\nu$ and their nonzero weights $\alpha_1,\ldots,\alpha_\nu$
from Lemma~\ref{lem:cartan-model}. Let $m$ be the nilpotence class of
$\mathfrak n$, and let $\mathcal Z$ be the finite list of all nonzero
bracket words of lengths at most $m$ in these generators. Repeated or
proportional vectors are allowed as different entries of the list. Each
word $Z$ has weight $\gamma_Z$, the sum of its letter weights with
multiplicity. The list spans $\mathfrak n$.

\subsection{Support exponents and normalized boxes}

For a nonempty compact set $E\subset W$, define
\begin{equation}\label{eq:support-exponents}
 s_i(E)=\max_{w\in E}\alpha_i(w),\qquad
 h_Z(E)=\sum_{X_i\text{ occurring in }Z}s_i(E),\qquad
 P(E)=\max_{\mathcal B}\sum_{Z\in\mathcal B}h_Z(E),
\end{equation}
where occurrences are counted with multiplicity and $\mathcal B$ runs
over the bases selected from $\mathcal Z$. The letters in a product may
be applied at different support positions. Their available amplitudes
therefore contribute the sum of the separate maxima $h_Z(E)$, which
can exceed $\max_E\gamma_Z$.
Fix an ordering of these
finitely many bases and choose the first maximizer whenever a choice is
needed. For $r\ge0$ and $L>0$, put
\begin{equation}\label{eq:weighted-box-definition}
 B_r(E,L)=\left\{\exp_N\left(\sum_{Z\in\mathcal B}
 t_Ze^{rh_Z(E)}Z\right): |t_Z|\le L\right\},
 \qquad B_r(E)=B_r(E,1).
\end{equation}
For example, let $\mathfrak n$ be the Heisenberg algebra with basis
$X,Y,Z$ and bracket $[X,Y]=Z$, and let a derivation $T$ act by
$[T,X]=X$, $[T,Y]=-Y$, and $[T,Z]=0$. Take the simply connected
group with Lie algebra $\mathfrak n\rtimes\R T$. Here $W=\R$, with
$\pi(T)=1$. With generators $X,Y$ and $E=[a,b]$, where $a\le b$,
\[
 s_X(E)=b,\qquad s_Y(E)=-a,\qquad h_Z(E)=b-a,
 \qquad P(E)=2(b-a).
\]
Thus the zero-weight direction $Z$ has width $e^{r(b-a)}$ in the box:
its width comes from the two supporting heights, although
$\gamma_Z=0$. This example also exhibits the Cartan overlap:
we may take $\mathfrak d=\Span(T,Z)$, so
$N\cap D=\exp(\R Z)$.

The next lemma controls the dependence on the maximizing basis choice.
A coordinate bound always refers to the displayed scaled basis, rather
than to an unscaled Euclidean ball.

\begin{lemma}[Box calculus]\label{lem:box-calculus}
There are constants depending only on the fixed Lie model with the
following properties, uniformly for $r\ge0$ and nonempty compact $E$.
\begin{enumerate}
\item The unit cube in a maximizing scaled basis and the Minkowski sum of the
intervals $[-e^{rh_Z(E)}Z,e^{rh_Z(E)}Z]$, $Z\in\mathcal Z$, contain
constant multiples of one another. Different maximizing bases therefore
give boxes comparable by fixed changes of $L$.
\item In these coordinates all Lie bracket coefficients are bounded
uniformly. Multiplication and inversion in exponential coordinates are
polynomials of uniformly bounded degree and coefficients. For some
$C,b$, a product of $k\ge1$ elements of $B_r(E,L)$ lies in
\[
 B_r(E,C(1+k)^b(1+L)^b).
\]
\item Writing $d_N=\dim N$, Haar measure satisfies
\begin{equation}\label{eq:box-volume}
 \vol_N B_r(E,L)\asymp L^{d_N}e^{rP(E)}.
\end{equation}
\item Exact weight dilations shift the template:
$S_{rw}B_r(E,L)$ is comparable, with fixed radius factors, to
$B_r(E+w,L)$. For each fixed $C_0$ there are $C,b$ such that, for
$r\ge1$ and $|d|_D\le C_0r$,
\begin{equation}\label{eq:box-conjugation}
 \Ad_d B_r(E,L)\subseteq
 B_r\left(E+\frac{\pi d}{r},C(1+r)^b(1+L)^b\right).
\end{equation}
The same assertion applies to $d^{-1}$.
\end{enumerate}
\end{lemma}

\begin{proof}
Write a word $Y\in\mathcal Z$ in a maximizing basis as
$Y=\sum_{Z\in\mathcal B}c_{Y,Z}Z$. Whenever $c_{Y,Z}\ne0$, replacing
$Z$ by $Y$ gives another basis. Maximality gives $h_Y(E)\le h_Z(E)$.
Thus
\begin{equation}\label{eq:basis-exchange-coefficients}
 e^{rh_Y(E)}Y=
 \sum_{Z\in\mathcal B}c_{Y,Z}e^{r(h_Y(E)-h_Z(E))}
 \big(e^{rh_Z(E)}Z\big)
\end{equation}
has uniformly bounded coefficients. There are only finitely many
unscaled bases and words, so their coefficients admit a common bound.
The reverse inclusion follows because the basis words belong to the
list. This proves the first assertion.

The bracket of two basis words is either zero or a word of the list,
with exponent equal to the sum of their exponents. Formula
\eqref{eq:basis-exchange-coefficients} therefore bounds every normalized
structure constant. The Baker--Campbell--Hausdorff formula terminates
at degree $m$, which proves the assertion for a fixed number of
factors. To control $k$ factors, expand the formal product of their
exponentials and then its logarithm, retaining only degrees at most
$m$. In degree $j$ there are at most $C_m(1+k)^j$ choices of factor
indices, and each term contains at most $j$ input coordinates.
Converting these finitely many homogeneous Lie polynomials into
brackets costs constants depending only on $m$. The asserted polynomial
bound in $k,L$ follows from the bounded normalized structure constants.

Haar measure of a simply connected nilpotent group in exponential
coordinates is a fixed multiple of Lebesgue measure. Indeed the
Jacobian factor of the exponential map is
$\det((1-e^{-\ad v})/\ad v)=1$, interpreted by its finite power
series, since $\ad v$ is nilpotent. The determinant of the scaled basis
is its unscaled determinant times $e^{rP(E)}$. The finitely many
unscaled determinants are bounded above and away from zero. This proves
\eqref{eq:box-volume}.

For an exact weight dilation,
$h_Z(E+w)=h_Z(E)+\gamma_Z(w)$. Every basis has total weight zero:
it is a weight-homogeneous basis of $\mathfrak n$, and the total
weight is the trace character, which vanishes by
Lemma~\ref{lem:canonical-height}. Hence maximizing bases remain maximizing
after a translation of $E$, and the claimed shift follows.

For conjugation, first remove the scalar action on each weight space.
Lemma~\ref{lem:canonical-height} bounds the remaining operator and its
inverse polynomially in $1+|d|_D$. Words of weight $\alpha_i$ span
the weight space containing $X_i$, and each such word satisfies
\begin{equation}\label{eq:weight-width-domination}
 h_Z(E)=\sum_{X_j\text{ in }Z}\max_E\alpha_j
 \ge\max_E\sum_{X_j\text{ in }Z}\alpha_j=s_i(E).
\end{equation}
Consequently the residual action on a scaled letter is a
polynomially bounded combination of scaled words. Applying brackets
and using the already established structure-constant bounds proves
the same assertion for every scaled word. Restoring the character
action is the exact shift by $\pi d/r$. This gives
\eqref{eq:box-conjugation}; replacing $d$ by $d^{-1}$ proves its
inverse version.
\end{proof}

Fix Euclidean norms on all the finite-dimensional spaces used below.
For sequences $r\to\infty$, a containment in $B_r(E,e^{o(r)})$
means that there is one sequence $\eta_r\to0$ for which all the
points under consideration lie in $B_r(E,e^{\eta_r r})$.
For a finite construction there is one common such envelope for all
its elements. This convention makes no assertion uniform over all
possible vanishing sequences. The box calculus shows that a fixed or
polynomial number of products preserves such a bound. Changes of
$E$ by Hausdorff distance $o(1)$ are also absorbed: each $h_Z$ is
Lipschitz in that distance, and the all-word description compares the
possibly different maximizing bases.

\begin{lemma}[Support symmetries]\label{lem:support-symmetries}
The function $P$ is translation invariant, depends only on the convex
hull, and is Lipschitz for Hausdorff distance. The function
\[
 p(w)=P([0,w])
\]
is a polyhedral norm on $W$. The group
\begin{equation}\label{eq:support-symmetry-group}
 \mathcal A=\{A\in\GL(W):P(AE)=P(E)\text{ for every nonempty
 compact convex }E\subset W\}
\end{equation}
is finite. Its pullback action permutes the oriented rays of the
generator weights $\alpha_i$.
\end{lemma}

When recording the dependence on the model in
Theorem~\ref{thm:bounded-height}, we write $\A_G=\mathcal A$ for
this same group.

\begin{proof}
Every basis from $\mathcal Z$ contains all the letters $X_i$:
all longer words vanish in $\mathfrak n/[\mathfrak n,\mathfrak n]$,
and the letters project to its specified basis. As observed above,
the sum of the word weights of any such basis is zero. Expanding this
identity in letters gives
\begin{equation}\label{eq:positive-weight-relation}
 \sum_i c_i\alpha_i=0,\qquad c_i\text{ positive integers}.
\end{equation}
The total exponent of a basis is unchanged by translating $E$, which
proves translation invariance. The other two assertions about $P$
follow from the corresponding properties of support functions and the
finiteness of the list.

For a segment $[0,w]$, each support is $\max(0,\alpha_i(w))$.
These expressions give positive homogeneity, piecewise linearity,
and the triangle inequality for $p$, the last by taking the maximum
of the basis sums after applying subadditivity term by term.
Translation invariance gives $p(-w)=p(w)$. If $w\ne0$, the weights
span $W^*$, and \eqref{eq:positive-weight-relation} implies that some
$\alpha_i(w)>0$. All word exponents on the segment are nonnegative,
and the mandatory letter $X_i$ makes every basis sum positive.
Thus $p$ is a norm.

Every member of $\mathcal A$ is a linear isometry of this norm.
Its unit ball is a full-dimensional bounded polytope, so its linear
isometry group injects into the permutation group of its finitely many
vertices. In particular $\mathcal A$ is finite.

Suppose $A\in\mathcal A$ and $A^*\alpha_j$ is not on one of the
original oriented rays. From a Euclidean ball $K$ centered at zero, remove
a sufficiently small cap at the support point of $A^*\alpha_j$.
The resulting compact convex set $E$ has the same supports in all
original directions $\alpha_i$, but a strictly smaller support in
direction $A^*\alpha_j$. Its supports in every other direction can
only decrease. Since each candidate basis sum for $P(AE)$ contains
the letter $j$, all these finitely many sums decrease strictly.
Thus $P(AE)<P(AK)=P(K)=P(E)$, contrary to $A\in\mathcal A$.
Pullback therefore maps the finite set of rays into itself; applying the
same argument to $A^{-1}$ proves that it permutes them.
\end{proof}

\subsection{Product schemes that remain nondegenerate}

A generator product scheme is a finite list of indices
$i_1,\ldots,i_k$, with endpoint map
\begin{equation}\label{eq:generator-endpoint-map}
 (t_1,\ldots,t_k)\longmapsto
 \prod_{j=1}^k\exp_N(t_je^{rs_{i_j}(E)}X_{i_j}).
\end{equation}
We use its logarithm in the maximizing scaled coordinates. This is
a polynomial map. Besides generating a box, later volume arguments
need a scheme whose normalized limiting map has full differential
rank somewhere. These are separate conclusions of the next lemma.

\begin{lemma}[Uniform product schemes]\label{lem:product-schemes}
For each fixed $L<\infty$ there are integers and constants depending
only on $L$ and the Lie model such that every element of $B_r(E,L)$,
for every $r\ge0$ and nonempty compact $E$, is the endpoint of a
scheme of at most that many factors, with all $|t_j|$ bounded by
that constant. Choices may be made Borel measurably in the endpoint.

For any sequence of these normalized coordinate systems there is a
subsequence and one fixed generator product scheme whose normalized
polynomial maps converge, with bounded degree and coefficients, to a
map submersive at an interior point of a fixed parameter cube. The
parameter point may be chosen to map to the identity.

If $r\to\infty$ and the radius is $e^{o(r)}$, generation still uses a
bounded number of factors, now with amplitudes
$e^{rs_i(E)+o(r)}$. Finally, every element of $N$ is a product of at
most a fixed number of the elements $\exp(tX_i)$, with unrestricted
real parameters.
\end{lemma}

\begin{proof}
We construct a submersive scheme for each limiting normalized group.
Its persistence under nearby group laws will then give uniform generation
by compactness.

Pass first to a subsequence with the same maximizing basis. The
normalized structure constants are bounded by
Lemma~\ref{lem:box-calculus}, so a further subsequence converges.
The Jacobi identity and the nilpotence-class bound pass to the
limit, and its group law is the limiting BCH polynomial.
Every letter is a basis vector. More generally, the bracket word
giving a basis vector satisfies the exact identity
\[
 Z(e^{rs_1(E)}X_1,\ldots,e^{rs_\nu(E)}X_\nu)=e^{rh_Z(E)}Z.
\]
Thus its normalized value is the same coordinate vector throughout
the sequence and in its limit. The limiting letters therefore still
bracket-generate the whole limiting algebra.

Here is a finite endpoint scheme with full rank in this limiting
group. Let $H$ be the subgroup generated by its letter
one-parameter subgroups, and consider
\[
 U=\Span\{\Ad_hX_i:h\in H,\ 1\le i\le\nu\}.
\]
This space contains the letters and is invariant under
$\Ad_{\exp(tX_i)}$. Differentiating proves its invariance under
$\ad X_i$, so it contains all brackets and is the entire algebra.
Choose $d_N$ vectors $\Ad_{h_j}X_{i_j}$ forming a basis, with each
$h_j$ a finite product of generator exponentials. The map
\begin{equation}\label{eq:submersive-conjugate-scheme}
 \Phi(t)=\prod_{j=1}^{d_N}h_j\exp(t_jX_{i_j})h_j^{-1}
\end{equation}
sends zero to the identity and has these basis vectors as its
differential columns at zero. Expanding the $h_j$ and their inverses
gives a generator scheme as in \eqref{eq:generator-endpoint-map}.
At a parameter point having fixed entries for the conjugating words
and zero entries for the inserted increments, its differential is
surjective. Enlarge a fixed parameter cube so this point is interior.

Use the same words and fixed parameters for the converging laws.
Their endpoint maps converge with derivatives on compact parameter
sets. The inverse function theorem, applied to the indicated
$d_N$ variables with the others fixed, supplies one fixed identity
neighborhood contained in their images for all sufficiently late
indices. Its local inverse supplies continuous parameter choices.
For a fixed-radius log cube choose $k$ so large that division of
every log vector by $k$ lies in this neighborhood. The identity
$\exp(v)=(\exp(v/k))^k$ covers the cube by a fixed number of copies
of the scheme.

This proves the uniform generation assertion by contradiction:
any sequence violating every proposed common length and parameter
bound has a convergent normalized-law subsequence to which the
preceding argument applies. Equivalently, the compact closure of
the normalized laws has a finite cover by neighborhoods carrying
the preceding inverse charts. There are only finitely many
maximizing bases. One may pad words by identities, and concatenate
the finitely many schemes, so there is no measurable choice among
infinitely many alphabets. Local inverse choices on finite Borel
partitions give the stated measurable endpoint choices.

For the growing-radius assertion let
$a_0=\min_i\|\alpha_i\|>0$. Enlarge $E$ to
$E+\varepsilon_r\overline B_W(1)$, where $\varepsilon_r=o(1)$ is
large enough that $r\varepsilon_ra_0$ dominates the logarithm of
the radius and the fixed box-comparison constants. Every word
exponent increases by at least $\varepsilon_ra_0$. The all-word
comparison then embeds the original growing box in a fixed-radius
box for the enlarged template. Generation there changes each letter
amplitude by only $e^{o(r)}$.

Finally choose one sufficiently large centered ball as template.
All its word exponents are strictly positive, so its boxes exhaust
$N$ as $r\to\infty$. The uniform number of factors just established
proves the last assertion.
\end{proof}

The measurable choices in this lemma do not imply any assertion
about the volume of their parameter sets. Whenever we apply a
polynomial image-volume estimate, we will instead identify an
independently sampled parameter domain, a fixed submersive limiting
scheme, and the measure of the successful parameters.

\subsection{Short switches and controlled paths}

At a supporting height, an exponentially large generator displacement
can be made by a short path in $G$. The next two lemmas quantify this
observation and assemble such switches into paths through prescribed
height sets.

\begin{lemma}[Logarithmic distortion]\label{lem:log-distortion}
There is a constant $C$ such that
\[
 |n|_G\le C\log(2+\|\log_Nn\|)\qquad(n\in N).
\]
For fixed constants controlling $|d|_D/r$ and the amplitude factor,
if $|d|_D=O(r)$,
$|t|\le e^{rs_i(E)+o(r)}$, and
$\alpha_i(\pi d)=rs_i(E)+o(r)$, then the distance between
$xnd$ and $xn\exp(tX_i)d$ is $o(r)$, uniformly in $x,n$.
With bounded amplitude factor and exact equality of heights, this
distance is $O(\log(2+r))$.
\end{lemma}

\begin{proof}
If $\|\log_Nn\|\le e^r$, choose a fixed large centered ball $E$
whose every word exponent is at least one. The finiteness of the
unscaled bases gives $n\in B_r(E,L_0)$ for a fixed $L_0$.
Lemma~\ref{lem:product-schemes} writes it as a bounded product of
letter factors with amplitudes at most $Me^{rs_i(E)}$.
Choose $v_i\in\mathfrak d$ with $\alpha_i(\pi v_i)=1$.
Conjugating a letter factor back by
$d_i=\exp_D((rs_i(E)+b\log(2+r))v_i)$ removes its exponential
amplitude. For a sufficiently large fixed $b$, the extra logarithmic
term also dominates the polynomial unipotent factor from
Lemma~\ref{lem:canonical-height}. The conjugated-back log vector is
bounded, while $|d_i|_G=O(r+1)$. The resulting outward, bounded,
and return moves prove the first estimate, also for bounded $n$.

By left invariance the distance for a switch is the length of
$\Ad_{d^{-1}}\exp(tX_i)$. Its logarithm has norm at most
\[
 C(1+r)^b|t|e^{-\alpha_i(\pi d)}.
\]
The first estimate makes this $o(r)$ in the stated regime and
$O(\log(2+r))$ in the exact regime.
\end{proof}

Write $a_D:D\to V=\mathfrak d/[\mathfrak d,\mathfrak d]$ for
abelianization, so that $\pi|_D$ factors through $a_D$.

\begin{lemma}[Paths visiting support positions]\label{lem:support-switch-paths}
Fix a compact convex $E_V\subset V$, a point $v_0\in E_V$, a
constant $L$, and pivots $d_0\in D$ satisfying
$a_D(d_0)=rv_0$ and $|d_0|_D\le C_0r$, where $r\ge1$.
Every $u\in B_r(\pi E_V,L)$ can be reached from $d_0$ at endpoint
$ud_0$ by a path with the following properties:
\begin{enumerate}
\item it uses a bounded number of $D$-legs and generator switches,
has length at most $Cr$, and all its $D$-positions have intrinsic
length at most $Cr$;
\item on each $D$-leg the abelianized position lies in $rE_V$;
each connecting switch has length at most $C\log(2+r)$;
\item for any fixed continuous homomorphism $l:G\to W$, all
path values of $l$ lie within $O_l(\log(2+r))$ of $r\,l(E_V)$,
where $l$ on $V$ denotes its induced linear map.
\end{enumerate}
Choices may be made measurably in $u$. The analogous statements
with switch length $o(r)$ hold for a subexponential box radius.
In particular, for a fixed compact convex $E\subset W$ and any
pivot of intrinsic length at most $C_0r$ with $\pi d_0\in rE$,
the same construction reaches $ud_0$, $u\in B_r(E,L)$, with
absolute heights in $rE+O(\log(2+r))$ and total length $O(r)$.
The constants in this specialization are independent of the
pivot's other abelianized coordinates.
One may also use a prescribed $D$-curve of length $O(r)$ and
intrinsic positions $O(r)$ whose height image attains all the
required supports, by visiting and reversing portions of that curve.
\end{lemma}

\begin{proof}
Choose, for each letter, a support point $v_i\in E_V$ for
$\alpha_i\circ\pi$. Fix a linear section $V\to\mathfrak d$.
From the pivot visit
$d_i=d_0\exp_D(r\widetilde{v_i-v_0})$, along the corresponding
one-parameter segment. Its abelianization traces the segment
$r[v_0,v_i]\subset rE_V$, and its length is $O(r)$.
Decompose $u$ using Lemma~\ref{lem:product-schemes}. To apply its
next letter factor, move along this leg on the current $D$-chart,
switch at $d_i$, and reverse the leg on the new sheet. The support
equality and Lemma~\ref{lem:log-distortion} give the claimed switch
cost. After the final return the endpoint is $ud_0$.

Every continuous homomorphism $l$ kills $N$, since its differential
kills $\mathfrak n\subset[\mathfrak g,\mathfrak g]$, and factors
on $D$ through $V$. It therefore has the specified values
on the legs. It is Lipschitz for the fixed metrics, so each short
switch changes its values by at most its Lipschitz constant times
the switch length. This proves the third assertion.

Measurable generator choices are supplied by
Lemma~\ref{lem:product-schemes}. Short connecting paths may also
be chosen measurably. The Riemannian exponential map on a closed
tangent ball of the asserted length radius is continuous from a
compact space and covers the relevant metric ball. Such a map has
a Borel section: use nested finite closed covers of successively
smaller diameters and choose the first subpiece whose image
contains the endpoint; all these image sets are compact, and the
nested choices have one limiting point. The corresponding
geodesics give the required choices. The same argument gives the
subexponential version.

For the specialization to $W$, choose an affine section of
$\pi:V\to W$ through $a_D(d_0)/r$ over $\pi d_0/r$, with one
fixed linear part, and apply the construction to the image of $E$.
All leg displacements depend only on differences in that affine
section. Their bounds therefore do not depend on its translating
kernel component; the bound on $|d_0|_D/r$ controls the positions.
For a prescribed curve, replace the displayed legs by its portions
between the pivot and a support position, and reverse each portion
after the switch. All estimates remain the same.
\end{proof}

\subsection{Packing after the multiplication map}

The upper and lower counting arguments have different sources.
For upper bounds we separate chosen lifts in $N\rtimes D$.
For lower bounds we remain on a fixed-$D$-position slice $Nd$, where the
parameterization by $N$ is injective. This distinction avoids
imposing a splitting hypothesis on $G$.

\begin{lemma}[Path packing]\label{lem:path-packing}
Fix $C_0>0$ and a separation distance $a>0$. There are constants
$C,b$ such that the following holds for $r\ge1$. Let $E\subset W$
be compact and contain zero. The endpoints of paths starting at
one point $x$, of length at most $C_0r$, and satisfying
\[
 \pi(x^{-1}\gamma(t))\in rE
 \quad\hbox{at every point of each path}
\]
have $a$-separated subsets of cardinality at most
\begin{equation}\label{eq:finite-path-packing}
 C(1+r)^b e^{rP(E)}.
\end{equation}
Each endpoint has a relative lift $(u,d)$ with
\[
 u\in B_r(E,C(1+r)^b),\qquad |d|_D\le Cr.
\]
These containments remain valid, after increasing $C,b$, under
fixed bounded right thickening in $N\rtimes D$.

In particular, along a sequence with relative path heights in
$r(E+o(1))$ and length $O(r)$, the upper separated count is
$e^{rP(E)+o(r)}$, and the lifts have $u\in B_r(E,e^{o(r)})$,
$|d|_D=O(r)$. Without a height restriction, a path of length $t$
has a relative lift with $|d|_D\le C(1+t)$ and
$\|\log_Nu\|\le e^{C(1+t)}$.
\end{lemma}

\begin{proof}
Translate the starting point to the identity and subdivide the
path into at most $C(1+r)$ increments of bounded length.
Lemma~\ref{lem:cartan-model} supplies bounded lifts $(v_j,e_j)$ of
these increments. Their accumulated product has coordinates
\[
 d_j=e_1\cdots e_j,
 \qquad
 u_j=\prod_{k=1}^j\Ad_{d_{k-1}}v_k.
\]
Thus $|d_j|_D\le C(1+r)$, and $\pi d_j$ is the actual height of
the corresponding path vertex.

Expand $\log_Nv_k$ in the fixed weight spaces. The polynomial
residual action gives, in a weight $\gamma$, coefficients bounded
by $C(1+r)^b e^{\gamma(\pi d_{k-1})}$. Words of that weight span
its space, and
$\max_E\gamma\le h_Z(E)$ for every such word $Z$.
Consequently each conjugated increment has normalized coordinates
bounded polynomially in $r$. The product estimate in
Lemma~\ref{lem:box-calculus} gives the asserted bound for $u_j$.

Bounded right thickening replaces $(u,d)$ by
$(u\Ad_dv,de)$ with $(v,e)$ bounded. The same estimate at the
terminal height controls this extra factor and proves the
thickening assertion.

Equip $N\rtimes D$ with a fixed left invariant Riemannian metric.
The homomorphism $q$ is globally Lipschitz: its derivative bound at
the identity is transported everywhere by left translations.
Therefore $a$-separated projected endpoints have separated chosen
lifts, with separation bounded below by a fixed positive number.
Small fixed-radius balls around those lifts are disjoint.
The thickening assertion places all of them in
\[
 B_r(E,C(1+r)^b)\times B_D(Cr+C).
\]
Product Haar is valid by Lemma~\ref{lem:canonical-height}; the first
factor has volume at most a polynomial times $e^{rP(E)}$, and the
second has polynomial volume by Lemma~\ref{lem:d-scaling}.
Every small lifted ball has the same positive Haar volume.
This proves \eqref{eq:finite-path-packing}. Injectivity of $q$ is
not used.

For the sequential statement, enlarge $E$ by a vanishing-radius
ball so it contains the actual relative path heights divided by
$r$ and zero. The support functional changes by $o(1)$, and the
box exponents by $o(1)$. Polynomial factors are subexponential.
For an unrestricted path of length $t$, its intermediate heights
are $O(1+t)$ because $\pi$ is Lipschitz. The same lifted-increment
argument now bounds each coordinate by $e^{O(1+t)}$, and its
polynomially many BCH terms preserve this bound.
\end{proof}

\begin{lemma}[Polynomial packing at bounded height]
\label{lem:polynomial-path-packing}
Fix $M,C_0,a>0$. There are $C,b$ such that, for every $t\ge0$,
the endpoints of paths starting at one point $x$, of length at
most $C_0t$, whose every relative height has norm at most $M$,
have $a$-separated subsets of cardinality at most
$C(1+t)^b$. The constants depend only on the displayed data and
the fixed Lie model. This lemma also holds when $N=1$.
\end{lemma}

\begin{proof}
If $N=1$, then $G=D$ is nilpotent. The endpoints lie in the
radius-$C_0t$ ball, and disjoint fixed-radius balls about separated
points, together with Lemma~\ref{lem:d-scaling}, give the polynomial
bound directly. Suppose now that $N\ne1$.
Repeat the bounded-increment lifting in the preceding proof.
Every accumulated $D$-coordinate has intrinsic length $O(1+t)$
and height norm at most $M$. On each of the finitely many weight
spaces the character factor is bounded by a constant depending
on $M$, while the residual action is polynomial in $1+t$.
Hence every conjugated $N$-increment has polynomial log norm.
There are $O(1+t)$ increments; the fixed-degree nilpotent product
formula bounds the log norm of their product by $C(1+t)^b$.
The same bounds survive fixed bounded right thickening.
The allowed $N$ log-coordinate region and the intrinsic $D$-ball
both have polynomial Haar volume. Disjoint small balls around
lifts of separated endpoints give the result exactly as before.
In particular this is a finite polynomial bound; it does not
follow by interpreting a subexponential error as polynomial.
\end{proof}

\begin{lemma}[Packing at a fixed $D$-position]\label{lem:fixed-sheet-packing}
For every $a>0$ there is $C_a<\infty$ such that, for every $d\in D$,
$x\in G$, and bounded measurable $S\subset N$, the set
$\{xnd:n\in S\}$ contains an $a$-separated subset with cardinality
at least $\vol_N(S)/C_a$. The constant is independent of $x,d,S$.
For a fixed finite number of such slices the corresponding assertion
holds for product Haar measure and the product maximum metric.
\end{lemma}

\begin{proof}
Set $K_a=N\cap\overline B_G(1,a)$. This is compact since $N$ is
closed and the metric on $G$ is proper. If $nd$ and $n_0d$ have
distance at most $a$, then
$n_0^{-1}n\in\Ad_dK_a$.
Conjugation by $d$ preserves $N$-Haar, so the possible $n$ have
measure at most $\vol_N K_a$, independently of $d$. A maximal
$a$-separated subset of the bounded slice covers it by radius-$a$
balls. Summing the preceding measure bound gives the assertion,
after a harmless increase of $a$ or of $C_a$ at boundary points.
Left translation by $x$ is an isometry. In the product maximum
metric the covering balls are products, so the same proof
multiplies the constants.
\end{proof}

The next estimate is needed when the $D$-positions vary. It
controls how much $N$-volume can be hidden in one projected ball,
including volume arising from $N\cap D$.

\begin{lemma}[Multiplicity of variable lifts]\label{lem:lift-multiplicity}
Fix a radius $R$. There are $C,b,c$ such that, for $T,H\ge0$ and
any metric ball $B_G(y,R)$, the set
\[
 \{u\in N:\text{ for some }d\in D, |d|_D\le T,
 \ \|\pi d\|\le H,\ q(u,d)\in B_G(y,R)\}
\]
has $N$-Haar outer measure at most
$C(1+T)^b e^{cH}$. In particular, if $T=O(r)$ and $H=o(r)$,
this measure is $e^{o(r)}$, with one common envelope for any
family subject to those bounds.
\end{lemma}

\begin{proof}
If the set is nonempty, fix one representative $(u_0,d_0)$.
Every other projected point differs from it by an element of a
fixed radius-$2R$ ball. Choose a bounded lift $(v,e)$ of that
difference, using Lemma~\ref{lem:cartan-model}. In $G$,
\[
 ud=u_0d_0ve=u_0\Ad_{d_0}(v)d_0e.
\]
It follows that
\begin{equation}\label{eq:overlap-multiplicity-formula}
 u_0^{-1}u=\Ad_{d_0}(v)a,
 \qquad a=d_0ed^{-1}\in N\cap D,
 \qquad |a|_D\le2T+O_R(1).
\end{equation}
The first factor has log norm at most
$C(1+T)^b e^{cH}$ by the character and unipotent bounds.
The second has polynomial log norm by intrinsic nilpotent
coordinates in $D$. The logarithm of $a$ in $D$ and in $N$ is
the same vector in $\mathfrak d\cap\mathfrak n$, because both
are restrictions of the injective exponential map of $G$.
The finite BCH formula therefore places $u_0^{-1}u$ in an
$N$ log-coordinate ball of radius
$C(1+T)^{b'}e^{c'H}$. Its Lebesgue, hence Haar, volume has the
stated form after changing the constants. This proves the bound
even for outer measure and uses no centrality of the overlap.
\end{proof}

We have obtained the required two directions of volume comparison.
Bracket-generated endpoints at a common $D$-position yield separated
points by Lemma~\ref{lem:fixed-sheet-packing}; images of paths with
controlled heights have bounded separated counts by
Lemma~\ref{lem:path-packing}. When a later construction first gives
$N$-volume at varying sublinear-height $D$-positions,
Lemma~\ref{lem:lift-multiplicity} supplies the remaining projection
bound.
